\documentclass[11pt,letterpaper]{article}
\usepackage[T1]{fontenc}
\usepackage[utf8]{inputenc}
\usepackage{newtxtext}
\usepackage[margin=1in,headheight=15pt]{geometry}
\usepackage{amsmath,amsthm,mathtools}
\usepackage{newtxmath}
\usepackage{microtype,booktabs,array,longtable}
\usepackage{xcolor}
\usepackage{enumitem,fancyhdr,titlesec}
\usepackage[most]{tcolorbox}
\usepackage{hyperref}
\definecolor{OliveInk}{HTML}{49523A}
\definecolor{WarmPaper}{HTML}{F4F3EA}
\definecolor{RuleGray}{HTML}{B7B9A8}
\hypersetup{colorlinks=true,linkcolor=OliveInk,citecolor=OliveInk,urlcolor=OliveInk,
 pdftitle={Oscillating Log-Concavity in Geometric Generating Functions},
 pdfauthor={Research manuscript prepared with ChatGPT},
 pdfsubject={Heim--Neuhauser Challenges 2 and 3; INVERT transform; Lambert series}}
\urlstyle{same}
\titleformat{\section}{\large\bfseries\color{OliveInk}}{\thesection}{0.7em}{}
\titleformat{\subsection}{\normalsize\bfseries\color{OliveInk}}{\thesubsection}{0.7em}{}
\setlength{\parindent}{0pt}
\setlength{\parskip}{5pt plus 1pt minus 1pt}
\setlength{\emergencystretch}{2em}
\linespread{1.035}
\pagestyle{fancy}
\fancyhf{}
\fancyhead[L]{\small\color{OliveInk}Oscillating log-concavity}
\fancyhead[R]{\small Research manuscript}
\fancyfoot[C]{\thepage}
\renewcommand{\headrulewidth}{0.3pt}
\newtheorem{theorem}{Theorem}[section]
\newtheorem{proposition}[theorem]{Proposition}
\newtheorem{lemma}[theorem]{Lemma}
\newtheorem{corollary}[theorem]{Corollary}
\theoremstyle{definition}
\newtheorem{definition}[theorem]{Definition}
\newtheorem{remark}[theorem]{Remark}
\newcommand{\N}{\mathbb N}
\newcommand{\Z}{\mathbb Z}
\newcommand{\R}{\mathbb R}
\newcommand{\C}{\mathbb C}
\newcommand{\Q}{\mathbb Q}
\newcommand{\Delt}{\Delta}
\newcommand{\sig}{\sigma}
\newcommand{\ii}{\mathrm i}
\DeclareMathOperator{\Rea}{Re}
\DeclareMathOperator{\Ima}{Im}
\DeclareMathOperator{\disc}{disc}
\DeclareMathOperator{\Gal}{Gal}
\newtcolorbox{summarybox}{enhanced,breakable,colback=WarmPaper,colframe=RuleGray,
 boxrule=0.5pt,arc=1.5pt,left=9pt,right=9pt,top=7pt,bottom=7pt}
\setlist{nosep,leftmargin=1.8em,topsep=4pt}
\numberwithin{equation}{section}

\title{\vspace{-1.5em}\color{OliveInk}\textbf{Oscillating Log-Concavity in\\Geometric Generating Functions}\\[0.6em]
\large Two Heim--Neuhauser challenges, a general pole criterion,\\
and the INVERT transform of squares}
\author{\normalsize A research manuscript prepared for Vladimir Reshetnikov}
\date{19 September 2026}

\begin{document}
\maketitle
\thispagestyle{empty}

\begin{abstract}
Heim and Neuhauser asked whether the geometric transforms of the square
sequence and of each divisor-sum sequence have infinitely many failures of
log-concavity. We prove both assertions, in stronger forms. A general
meromorphic criterion shows that, after a unique positive dominant pole,
a first subdominant circle containing no positive real pole forces both
signs of the adjacent Tur\'an determinant to occur with bounded gaps.
For divisor weights, a parity identity and elementary radial asymptotics
produce a negative pole inside the unit disk. This proves the conclusion
for every real divisor exponent and every positive composition weight.
For integer power weights, Eulerian-polynomial reciprocity gives a complete
qualitative classification: exponents zero and one are eventually
log-linear and strictly log-concave, respectively, whereas every exponent
at least two exhibits both signs with bounded gaps.

For square weights, the determinant itself satisfies a third-order integer
recurrence. We prove that it never vanishes, that every four consecutive
indices contain both signs, and that each sign has density one-half within
every arithmetic progression. The density proof uses an irrational root
angle certified by the Galois group of a cubic of discriminant $-107$.
The archive includes a dependency-free exact-arithmetic verification program,
100,000 computed signs, and rational certificates for sample Lambert-series
inequalities. No finite computation is used to justify an infinite conclusion.
\end{abstract}

\begin{summarybox}
\textbf{Proven statements and bibliographical status.}
The targets are precisely Challenges 2 and 3 in Section 3 of
\cite{HN23}, arXiv:2302.13327v1, dated 26 February 2023.
That version and its challenge page were inspected directly. A targeted
search on 19 September 2026 did not locate a prior resolution; this is not
an exhaustive certification of priority. The manuscript supplies complete
arguments for the stated assertions, but has not undergone independent
peer review or formal proof-assistant verification. The known generating
function and recurrence of OEIS A033453 are explicitly credited rather
than presented as new.
\end{summarybox}

\clearpage
\tableofcontents
\clearpage

\section{The questions and the resulting theorems}
\label{sec:questions}

For a positive sequence $(a_n)_{n\ge0}$, write
\begin{equation}\label{eq:delta}
\Delt_n=a_n^2-a_{n-1}a_{n+1}\qquad(n\ge1).
\end{equation}
Thus $\Delt_n>0$ means strict log-concavity at $n$, while $\Delt_n<0$
means strict log-convexity there. Following the terminology of
Heim and Neuhauser, a negative determinant is an \emph{exception} to
log-concavity. Their two geometric families are
\begin{align}
\sum_{n\ge0}q^{\psi_d}(n)z^n
 &=\frac{1}{1-\sum_{k\ge1}k^d z^k},\label{eq:original-power}\\
\sum_{n\ge0}q^{\sig_d}(n)z^n
 &=\frac{1}{1-\sum_{k\ge1}\sig_d(k)z^k},
 &\sig_d(k)&=\sum_{\ell\mid k}\ell^d.\label{eq:original-divisor}
\end{align}
The constant coefficient in both definitions is $1$.

Challenge 2 of \cite{HN23} predicts infinitely many exceptions and
non-exceptions for $q^{\psi_2}$. Challenge 3 predicts infinitely many
exceptions for $q^{\sig_d}$, for every integer $d\ge0$.
These are questions with $d$ fixed and $n$ tending to infinity. They differ
from the main large-$d$, fixed-$n$ results in that paper. We do not address
the separate exponential/product-generating-function Challenge 1.

It is useful to retain an additional positive parameter $u$. Set
\begin{align}
A^\psi_{d,u}(z)&=\frac{1}{1-u\sum_{k\ge1}k^d z^k},
 &a^\psi_{d,u}(n)&=[z^n]A^\psi_{d,u}(z),\label{eq:power-family}\\
A^\sig_{s,u}(z)&=\frac{1}{1-u\sum_{k\ge1}\sig_s(k)z^k},
 &a^\sig_{s,u}(n)&=[z^n]A^\sig_{s,u}(z),\label{eq:divisor-family}
\end{align}
where $d$ is a nonnegative integer, $s$ may be any real number, and $u>0$.
The use of a positive parameter in this type of geometric transform is
already part of the framework in \cite{HN21}. For integer $s=d\ge0$ and
integer $u>0$, these are integer sequences with a colored-composition
interpretation. Allowing real $s$ or $u$ gives positive real sequences.

\begin{definition}
A set $S\subseteq\N$ has \emph{eventually bounded gaps} if there are integers
$N,L\ge1$ such that each block $\{n,n+1,\ldots,n+L-1\}$ with $n\ge N$
contains an element of $S$.
\end{definition}
This is stronger than infinitude, and it implies positive lower natural
density. The constants $N,L$ below are allowed to depend on the fixed
parameters; no uniformity in $d,s,u$ is asserted.

\begin{theorem}[Divisor weights]\label{thm:divisor-main}
For every $s\in\R$ and $u>0$, the sequence $a^\sig_{s,u}(n)$ has strict
log-concavity and strict log-convexity at sets of indices that both have
eventually bounded gaps. In particular, Challenge 3 of \cite{HN23} holds
for every $d\ge0$, and infinitely many strict non-exceptions occur as well.
\end{theorem}

\begin{theorem}[Classification for integer power weights]\label{thm:power-main}
For $u>0$, let $\Delt_n$ be the determinant of $a^\psi_{d,u}(n)$.
Then
\begin{equation}\label{eq:classification}
\begin{array}{c|c|c}
 d & \Delt_1 & \text{behavior for }n\ge2\\ \hline
 0 & -u & \Delt_n=0\\
 1 & -2u & \Delt_n=u^2\\
 d\ge2 & -u\,2^d & \text{both strict signs have eventually bounded gaps.}
\end{array}
\end{equation}
Consequently, Challenge 2 holds, and its infinitude conclusion extends
to every integer power exponent $d\ge2$ and every $u>0$.
\end{theorem}
The entry $\Delt_1=-u2^d$ follows directly from
$a_0=1$, $a_1=u$, and $a_2=u^2+u2^d$.

\begin{theorem}[Sharp square-weighted conclusions]\label{thm:square-main}
Let $a_n=a^\psi_{2,1}(n)$ and define $\Delt_n$ by \eqref{eq:delta}.
Then:
\begin{enumerate}[label=(\roman*)]
\item $\Delt_1=-4$, and $\Delt_n$ is odd for every $n\ge2$.
In particular, equality never occurs.
\item Every four consecutive positive indices contain both a positive
and a negative determinant. Runs of three of either sign do occur, so
this four-index block length is optimal.
\item For every modulus $m\ge1$ and residue class $r\pmod m$, each strict
sign has relative natural density $1/2$ among the positive indices in
that class. Equivalently,
\begin{equation}\label{eq:AP-density}
\lim_{X\to\infty}\frac{1}{X}
 \#\{1\le n\le X:n\equiv r\pmod m,\ \Delt_n>0\}
 =\frac{1}{2m},
\end{equation}
and the same limit holds with $\Delt_n<0$.
\end{enumerate}
The sign sequence is not eventually periodic.
\end{theorem}

Theorems \ref{thm:divisor-main} and \ref{thm:power-main} are proved in
Sections \ref{sec:lambert} and \ref{sec:power}, using the general analytic
criterion of Section \ref{sec:criterion}. The stronger square-case
statements occupy Section \ref{sec:squares}.

\section{Geometric transforms and the dominant pole}
\label{sec:framework}

Let
\begin{equation}\label{eq:general-transform}
G(z)=\sum_{k\ge1}g_kz^k,\qquad
A(z)=\frac{1}{1-uG(z)}=\sum_{n\ge0}a_nz^n,
\end{equation}
where all $g_k>0$ and $u>0$. Comparing coefficients in
$(1-uG)A=1$ gives the exact convolution recurrence
\begin{equation}\label{eq:convolution}
a_0=1,\qquad a_n=u\sum_{k=1}^n g_ka_{n-k}\quad(n\ge1).
\end{equation}
Alternatively, expanding $A=\sum_{j\ge0}u^jG^j$ gives
\begin{equation}\label{eq:compositions}
a_n=\sum_{\substack{k_1+\cdots+k_j=n\\j\ge1,\ k_i\ge1}}
 u^j\prod_{i=1}^j g_{k_i}\qquad(n\ge1).
\end{equation}
For integer weights, a part of size $k$ has $g_k$ possible colors, with
an additional $u$ choices per part. Formula \eqref{eq:compositions} also
makes positivity transparent.

\begin{lemma}[Unique positive dominant pole]\label{lem:dominant}
Suppose $G$ is analytic in $|z|<b$ and $G(x)\to+\infty$ as $x\uparrow b$,
where $0<b<\infty$. There is a unique $\rho\in(0,b)$ with
$uG(\rho)=1$. The function $A$ is meromorphic in $|z|<b$, and its only
pole in $|z|\le\rho$ is the simple pole at $\rho$. Its principal part can
be written
\begin{equation}\label{eq:C}
\frac{C}{1-z/\rho},\qquad C=\frac{1}{u\rho G'(\rho)}>0.
\end{equation}
In particular, for some $r>\rho$,
\begin{equation}\label{eq:dominant-asymptotic}
a_n=C\rho^{-n}+O(r^{-n}),\qquad
\rho^na_n\longrightarrow C,\qquad
\frac{a_n}{a_{n-1}}\longrightarrow\rho^{-1}.
\end{equation}
\end{lemma}

\begin{proof}
Strict positivity of the coefficients makes $G(x)$ strictly increasing
on $(0,b)$, with $G(0)=0$ and the stated divergent limit. This gives the
unique $\rho$ and $G'(\rho)>0$.

If $|z|<\rho$, then $|uG(z)|\le uG(|z|)<1$, so there is no pole.
If $|z|=\rho$ and $uG(z)=1$, equality holds in the triangle inequality
for $\sum_{k\ge1}g_kz^k$. All its nonzero summands must have the same
argument. In particular, the terms in $z$ and $z^2$ have the same argument,
so $z/\rho=1$. Thus the dominant circle has just the positive pole.
Expanding $1-uG(z)$ at $z=\rho$ gives \eqref{eq:C}.

After subtracting \eqref{eq:C}, the pole is removed. Zeros of a nonzero
analytic function are isolated; hence the difference is analytic on a
disk of radius larger than $\rho$. Cauchy's coefficient estimate gives
\eqref{eq:dominant-asymptotic}.
\end{proof}

The leading term alone does not decide log-concavity. A geometric
sequence has identically zero determinant:
\begin{equation}\label{eq:leading-cancellation}
(C\rho^{-n})^2-(C\rho^{-(n-1)})(C\rho^{-(n+1)})=0.
\end{equation}
Consequently, an exponentially small correction to $a_n$ can determine
the entire sign pattern of $\Delt_n$. The decisive information is often
the first \emph{subdominant} pole or group of poles, rather than a sharper
estimate of the positive leading constant.

\section{A general pole criterion for both signs}
\label{sec:criterion}

The next result is the common mechanism behind the two challenges.
We first give a short qualitative argument, then a quantitative-in-form
argument proving bounded gaps. Neither requires the subdominant poles
to be simple or to form a single conjugate pair.

\subsection{A short proof of infinitude using coefficient positivity}

We use the following elementary form of Pringsheim's principle.

\begin{lemma}[Positive coefficients force a positive singularity]
\label{lem:pringsheim}
A power series with eventually nonnegative real coefficients and finite,
positive radius of convergence $R$ is singular at the positive real point $R$.
\end{lemma}
\begin{proof}
Subtracting a polynomial reduces to nonnegative coefficients
$F(z)=\sum_{n\ge0}f_nz^n$. Suppose $F$ were analytic at $R$, with a
continuation to $|z-R|<\delta$. Choose $\varepsilon>0$ so small that
$3\varepsilon<\delta$ and $r=R-\varepsilon>0$.
The Taylor series of the continuation at $r$ converges at
$r+2\varepsilon=R+\varepsilon$. Its derivatives at $r$ equal the derivatives
of the original series. Nonnegativity permits interchange of the sums:
\begin{align*}
\sum_{k\ge0}\frac{F^{(k)}(r)}{k!}(2\varepsilon)^k
&=\sum_{n\ge0}f_n\sum_{k=0}^n\binom nk r^{n-k}(2\varepsilon)^k\\
&=\sum_{n\ge0}f_n(R+\varepsilon)^n<\infty.
\end{align*}
This contradicts the radius of convergence.
\end{proof}

\begin{proposition}[Qualitative pole obstruction]\label{prop:qualitative}
Let $A(z)=\sum a_nz^n$ have positive real coefficients and a unique
smallest-modulus pole $\rho>0$, simple, with principal part
$C/(1-z/\rho)$ and $C>0$. Suppose the remainder
\begin{equation}\label{eq:H}
H(z)=A(z)-\frac{C}{1-z/\rho}
\end{equation}
has finite radius of convergence $R>\rho$ but is analytic at the positive
point $R$. Then $\Delt_n$ is strictly positive infinitely often and strictly
negative infinitely often.
\end{proposition}
\begin{proof}
The pole expansion gives \eqref{eq:dominant-asymptotic}. If the determinants
were eventually nonnegative, the ratios $a_n/a_{n-1}$ would eventually
be decreasing to $\rho^{-1}$. Therefore $b_n=\rho^na_n$ would eventually
be increasing to $C$, so $b_n\le C$ eventually. The coefficients of
$-H$ would consequently be eventually nonnegative. Lemma
\ref{lem:pringsheim} would force a singularity at $R$, a contradiction.

If the determinants were eventually nonpositive, the ratios would be
increasing to $\rho^{-1}$, and $b_n$ would be decreasing to $C$.
Then $H$, rather than $-H$, would have eventually nonnegative coefficients,
giving the same contradiction. These two exclusions prove both strict
infinitude statements.
\end{proof}

This proof already resolves the infinitude questions once a suitable
secondary pole is found. The next theorem extracts more from the same
analytic situation.

\subsection{A bounded-gap theorem}

\begin{theorem}[Subdominant poles force both signs with bounded gaps]
\label{thm:spectral}
Let $A(z)=\sum_{n\ge0}a_nz^n$ be real at real regular points and meromorphic
in $|z|<B$. Suppose:
\begin{enumerate}[label=(\alph*)]
\item its unique pole of smallest modulus is a simple pole $\rho>0$,
with principal part $C/(1-z/\rho)$, where $C>0$;
\item there is at least one other pole in $|z|<B$;
\item if $R>\rho$ is the smallest modulus of these other poles, then
$A$ has no pole at the positive real point $R$.
\end{enumerate}
Then there are $N,L\ge1$ such that every block of $L$ indices beginning
at $n\ge N$ contains both a strict positive and a strict negative value
of $a_n^2-a_{n-1}a_{n+1}$.
\end{theorem}

\begin{proof}
\textbf{Step 1: isolate the first subdominant terms.}
There are finitely many poles on $|z|=R$. Let $m$ be their greatest
order, and list the poles of this order as
$r_j=Re^{\ii\theta_j}$, $1\le j\le J$. There is no $\theta_j$ congruent
to zero modulo $2\pi$, by hypothesis (c).
Subtracting all principal parts on and inside this circle leaves a
function analytic on a strictly larger disk. The coefficient identity
\[
[z^n](1-z/r)^{-k}=\binom{n+k-1}{k-1}r^{-n}
\]
therefore yields
\begin{equation}\label{eq:subdominant}
a_n=C\rho^{-n}+n^{m-1}R^{-n}\bigl(T_n+o(1)\bigr),
\qquad T_n=\sum_{j=1}^J c_je^{-\ii n\theta_j}.
\end{equation}
Here each $c_j$ is nonzero, and conjugate poles have conjugate constants,
so $T_n$ is real. The $o(1)$ consists of lower powers of $n$ and
exponentially smaller terms. In particular it remains $o(1)$ after
shifting $n$ by either $1$ or $-1$.

\textbf{Step 2: compute the first surviving determinant term.}
Put $a_n=C\rho^{-n}+e_n$. Direct expansion gives
\begin{equation}\label{eq:delta-cross}
\Delt_n=C\rho^{-n}
 \bigl(2e_n-\rho e_{n+1}-\rho^{-1}e_{n-1}\bigr)
 +e_n^2-e_{n-1}e_{n+1}.
\end{equation}
Substituting \eqref{eq:subdominant}, and using
$(n\pm1)^{m-1}/n^{m-1}=1+O(1/n)$, gives
\begin{equation}\label{eq:delta-asymptotic}
\Delt_n=C(\rho R)^{-n}n^{m-1}\bigl(U_n+o(1)\bigr),
\end{equation}
where
\begin{equation}\label{eq:U}
U_n=\sum_{j=1}^J v_je^{-\ii n\theta_j},
\qquad
v_j=c_j\left(2-\frac{\rho}{r_j}-\frac{r_j}{\rho}\right).
\end{equation}
The quadratic remainder in \eqref{eq:delta-cross}, divided by the
positive scale in \eqref{eq:delta-asymptotic}, is
$O(n^{m-1}(\rho/R)^n)$, which tends to zero.
No $v_j$ vanishes: $2-x-x^{-1}=0$ implies $x=1$, whereas
$r_j\ne\rho$.

\textbf{Step 3: a nonzero finite oscillation cannot keep one sign on
arbitrarily long blocks.}
Write $\omega_j=e^{-\ii\theta_j}$. The frequencies are distinct,
$|\omega_j|=1$, and none equals $1$. For every fixed $\omega\ne1$,
\begin{equation}\label{eq:geometric-average}
\left|\frac1L\sum_{n=M}^{M+L-1}\omega^n\right|
\le\frac{2}{L|1-\omega|},
\end{equation}
uniformly in the starting index $M$. Hence the mean of $U_n$ over such
blocks tends to zero uniformly in $M$.
Because $U_n$ is real and the list is closed under conjugation, the
constant term in $U_n^2$ is
\begin{equation}\label{eq:variance}
V=\sum_{j=1}^J|v_j|^2>0.
\end{equation}
Every other frequency in the square again has uniformly vanishing
block average by \eqref{eq:geometric-average}. Thus
\begin{equation}\label{eq:uniform-moments}
\frac1L\sum_{n=M}^{M+L-1}U_n\longrightarrow0,
\qquad
\frac1L\sum_{n=M}^{M+L-1}U_n^2\longrightarrow V,
\end{equation}
uniformly in $M$. Real poles at $-R$ cause no difficulty: their frequency
is $-1$, and its square correctly contributes to $V$.

Let $K=\sum_j|v_j|$, so $|U_n|\le K$, and choose
$\eta=V/(8K)>0$. Fix $L$ large enough that every length-$L$ block has
mean $U$ of absolute value at most $\eta$ and mean $U^2$ at least $V/2$.
If all values on one block were at most $\eta$, their positive parts
would have mean at most $\eta$. Consequently
\[
\operatorname{mean}(U^2)
\le K\operatorname{mean}|U|
=K\bigl(2\operatorname{mean}U_+-\operatorname{mean}U\bigr)
\le3K\eta=3V/8,
\]
a contradiction. Applying the same argument to $-U$ shows that every
such block contains a value above $\eta$ and one below $-\eta$.

Finally choose $N$ so that the error in \eqref{eq:delta-asymptotic} has
absolute value below $\eta/2$ for $n\ge N$. Its prefactor is positive,
so the two values in each block retain their respective signs in
$\Delt_n$.
\end{proof}

\begin{remark}[What the hypotheses do and do not say]
The theorem does not require that the nearest subdominant pole be negative
real. A nearer complex pole is equally useful. It does require exclusion
of a \emph{positive} pole on the first subdominant circle: such a pole
would contribute a nonoscillating term to $U_n$, potentially dominating
the sign. The power-weighted case $d=1$ below is a concrete example of
that different behavior.
\end{remark}

\section{Divisor weights: a negative pole for every real exponent}
\label{sec:lambert}

Put
\begin{equation}\label{eq:Lambert}
L_s(z)=\sum_{n\ge1}\sig_s(n)z^n
      =\sum_{k\ge1}\frac{k^s z^k}{1-z^k},\qquad |z|<1.
\end{equation}
The equality follows by expanding each geometric denominator and grouping
terms with the same product of indices. Absolute convergence is justified
by the polynomial bound
$\sig_s(n)\le n^{1+\max(s,0)}$.
Since $\sig_s(n)\ge1$, $L_s(x)\ge x/(1-x)$ for $0<x<1$.
Lemma \ref{lem:dominant} therefore applies to
$A^\sig_{s,u}=1/(1-uL_s)$ with $b=1$.

It remains to find a second pole before reaching $|z|=1$.
A negative real pole will suffice. The central observation is that,
despite the alternating signs of its Taylor expansion on $(-1,0)$,
$L_s(z)$ tends to \emph{positive} infinity as $z\to-1$ radially.

\subsection{An exact parity identity}

\begin{lemma}\label{lem:parity}
For real $s$, with $\sig_s(n/2)$ interpreted as zero when $n$ is odd,
\begin{equation}\label{eq:sigma-even}
\sig_s(2n)=(1+2^s)\sig_s(n)-2^s\sig_s(n/2).
\end{equation}
Consequently,
\begin{equation}\label{eq:Lambert-parity}
L_s(-z)=-L_s(z)+(2+2^{s+1})L_s(z^2)-2^{s+1}L_s(z^4).
\end{equation}
\end{lemma}
\begin{proof}
Write $n=2^a m$ with $m$ odd, and put $q=2^s$.
Then $\sig_s(n)=(1+q+\cdots+q^a)\sig_s(m)$.
The identity of finite geometric sums
\[
1+q+\cdots+q^{a+1}
=(1+q)(1+q+\cdots+q^a)-q(1+q+\cdots+q^{a-1})
\]
proves \eqref{eq:sigma-even}, including $a=0$ when the final sum is empty.
This argument also covers $s=0$, so no division by $2^s-1$ is needed.

The even-indexed part of $L_s(z)$ is
$(1+2^s)L_s(z^2)-2^sL_s(z^4)$ by \eqref{eq:sigma-even}.
Twice this part minus the whole series is $L_s(-z)$, yielding
\eqref{eq:Lambert-parity}.
\end{proof}

\subsection{Positive radial asymptotics}

\begin{lemma}\label{lem:positive-radial}
As $t\downarrow0$,
\begin{equation}\label{eq:positive-radial}
L_s(e^{-t})=
\begin{cases}
\Gamma(s+1)\zeta(s+1)t^{-s-1}(1+o(1)),&s>0,\\
t^{-1}\log(1/t)+O(t^{-1}),&s=0,\\
\zeta(1-s)t^{-1}(1+o(1)),&s<0.
\end{cases}
\end{equation}
\end{lemma}
\begin{proof}
For $s>0$, write
\[
t^{s+1}L_s(e^{-t})
=t\sum_{k\ge1}\frac{(kt)^s}{e^{kt}-1}.
\]
This is a Riemann sum for the integrable function
$f_s(x)=x^s/(e^x-1)$ on $(0,\infty)$.
For completeness, near zero one has $f_s(x)\le x^{s-1}$ and
$\sum_{k\le M}k^{s-1}=O_s(M^s)$; hence the portion $kt\le\varepsilon$
is $O_s(\varepsilon^s)$ uniformly as $t\downarrow0$.
The far tail is uniformly controlled by exponential decay, and the
remaining compact interval is handled by ordinary Riemann sums.
Finally, nonnegative-term integration gives
\[
\int_0^\infty\frac{x^s}{e^x-1}\,dx
=\sum_{j\ge1}\int_0^\infty x^se^{-jx}\,dx
=\Gamma(s+1)\zeta(s+1).
\]

For $s=0$, take $M=\lfloor1/t\rfloor$. Uniformly for $0<x\le1$,
$(e^x-1)^{-1}=x^{-1}+O(1)$, so
\[
t\sum_{k\le M}\frac{1}{e^{kt}-1}
=\sum_{k\le M}\frac1k+O(tM)=\log(1/t)+O(1).
\]
For $k>M$, the bound
$(e^{kt}-1)^{-1}\le e^{-kt}/(1-e^{-1})$ gives a contribution $O(1)$
to $tL_0(e^{-t})$. This proves the middle case.

For $s<0$, use
\[
tL_s(e^{-t})
=\sum_{k\ge1}k^{s-1}\frac{kt}{e^{kt}-1}.
\]
The factor on the right lies in $(0,1)$ and tends to $1$ for each fixed
$k$. Since $\sum k^{s-1}=\zeta(1-s)<\infty$, dominated convergence proves
the last case.
\end{proof}

\subsection{The negative radial limit and the pole}

\begin{proposition}\label{prop:negative-radial}
For every real $s$,
\begin{equation}\label{eq:negative-radial}
L_s(-e^{-t})=
\begin{cases}
2^{-s-1}\Gamma(s+1)\zeta(s+1)t^{-s-1}(1+o(1)),&s>0,\\
\frac12t^{-1}\log(1/t)+O(t^{-1}),&s=0,\\
2^{s-1}\zeta(1-s)t^{-1}(1+o(1)),&s<0.
\end{cases}
\end{equation}
In particular, $L_s(-r)\to+\infty$ as $r\uparrow1$.
\end{proposition}
\begin{proof}
Insert $z=e^{-t}$ into \eqref{eq:Lambert-parity} and apply
\eqref{eq:positive-radial} at $t,2t,4t$.
When $s>0$, the coefficient of the leading term is
\[
-1+(2+2^{s+1})2^{-s-1}-2^{s+1}4^{-s-1}=2^{-s-1}.
\]
When $s<0$, it is
\[
-1+\tfrac12(2+2^{s+1})-\tfrac14 2^{s+1}=2^{s-1}.
\]
For $s=0$, the coefficient of $t^{-1}\log(1/t)$ is
$-1+4/2-2/4=1/2$; replacing $\log(1/t)$ by
$\log(1/(2t))$ or $\log(1/(4t))$ only changes the $O(t^{-1})$ term.
All displayed leading constants are strictly positive.
\end{proof}

\begin{proof}[Proof of Theorem \ref{thm:divisor-main}]
Fix $s\in\R$ and $u>0$. The function $L_s(-r)$ is continuous for
$0\le r<1$, starts at zero, and tends to positive infinity.
Thus some $\beta\in(0,1)$ satisfies $uL_s(-\beta)=1$.
Since the numerator of $A^\sig_{s,u}$ is $1$, this zero of its denominator
is a genuine pole, of whatever finite order it may have.

Lemma \ref{lem:dominant} supplies the unique dominant pole $\rho>0$,
and also gives $\rho<\beta$. If $R$ is the smallest modulus of a
nondominant pole, then $\rho<R\le\beta<1$.
There is no positive pole at $R$, because $L_s(x)$ is strictly increasing
on $(0,1)$ and crosses $1/u$ exactly once.
Theorem \ref{thm:spectral}, with $B=1$, now proves the result.
\end{proof}

\begin{remark}
We did not prove that $-\beta$ is the nearest nondominant pole, or even
that the negative solution is unique. Neither assertion is needed.
This avoids a potentially difficult global zero-location problem for
Lambert series. Also, the argument does not prove the stronger empirical
odd/even pattern reported for small nonnegative integer exponents in
\cite{HN23}; it proves the stated infinitude challenge and the bounded-gap
strengthening without relying on that pattern.
\end{remark}

\section{Power weights: a complete integer-exponent classification}
\label{sec:power}

For an integer $d\ge1$, define the Eulerian polynomial $E_d(z)$ by
\begin{equation}\label{eq:Eulerian-GF}
\sum_{n\ge1}n^d z^n=\frac{zE_d(z)}{(1-z)^{d+1}}.
\end{equation}
We need only a few elementary properties, which we record with a proof.

\begin{lemma}\label{lem:Eulerian}
The polynomial $E_d$ has degree $d-1$, strictly positive coefficients,
$E_d(1)=d!$, and
\begin{equation}\label{eq:Eulerian-reciprocal}
E_d(z)=z^{d-1}E_d(1/z).
\end{equation}
Moreover,
\begin{equation}\label{eq:Eulerian-rec}
E_1(z)=1,\qquad
E_{d+1}(z)=(1+dz)E_d(z)+z(1-z)E'_d(z).
\end{equation}
\end{lemma}
\begin{proof}
The case $d=1$ is $\sum n z^n=z/(1-z)^2$.
Applying $z\,d/dz$ to \eqref{eq:Eulerian-GF} proves
\eqref{eq:Eulerian-rec} and the next generating function.
If $E_d(z)=\sum_{j=0}^{d-1}e_{d,j}z^j$, its coefficient recurrence is
\[
e_{d+1,j}=(j+1)e_{d,j}+(d-j+1)e_{d,j-1},
\]
with out-of-range entries zero. Starting from $e_{1,0}=1$, this proves
positivity, the degree, and symmetry under $j\mapsto d-j$ for $E_{d+1}$
by induction. The symmetry is \eqref{eq:Eulerian-reciprocal}.
Putting $z=1$ in \eqref{eq:Eulerian-rec} gives
$E_{d+1}(1)=(d+1)E_d(1)$, so $E_d(1)=d!$.
\end{proof}

The geometric transform becomes the rational function
\begin{equation}\label{eq:power-rational}
A^\psi_{d,u}(z)=\frac{(1-z)^{d+1}}{D_{d,u}(z)},\qquad
D_{d,u}(z)=(1-z)^{d+1}-uzE_d(z).
\end{equation}
The denominator has degree $d+1$, constant coefficient $1$, and leading
coefficient $(-1)^{d+1}$. Its value at $1$ is $-ud!\ne0$.
The numerator therefore shares no zero with the denominator: every root
of $D_{d,u}$ is a genuine pole. In particular, $z=1$ is a zero of the
rational function, not an extra pole.

There is a unique positive pole $\rho\in(0,1)$, and it is the unique
smallest-modulus pole, by Lemma \ref{lem:dominant} applied inside the
unit disk. To invoke Theorem \ref{thm:spectral}, we inspect possible
positive poles outside it.

\subsection{Even exponents at least two}

For even $d\ge2$ and real $x>1$, the numerator $xE_d(x)$ is positive
and $(1-x)^{d+1}$ is negative. Thus the rational continuation of
$\sum n^d z^n$ is negative at $x$, and
\[
1-u\frac{xE_d(x)}{(1-x)^{d+1}}>0.
\]
There is no positive pole beyond $1$. Since $D_{d,u}$ has degree
$d+1\ge3$ and the positive pole $\rho$ is simple, other poles exist,
but none of them is positive. Theorem \ref{thm:spectral} applies to this
rational function in a disk large enough to include its first
nondominant circle.

\subsection{Odd exponents at least three}

For odd $d$, $d+1$ is even and \eqref{eq:Eulerian-reciprocal} gives
\begin{equation}\label{eq:D-reciprocal}
D_{d,u}(z)=z^{d+1}D_{d,u}(1/z).
\end{equation}
The roots are therefore closed under inversion, with multiplicities.
The only positive root below $1$ is $\rho$, so the only positive root
above $1$ is $1/\rho$; there is no root at $1$.
Remove these two simple roots. When $d\ge3$, at least two roots remain,
and their multiset is still closed under inversion. At least one of them
has modulus at most $1$; otherwise its reciprocal would violate the same
claim. Consequently the nearest nondominant modulus satisfies
\[
\rho<R\le1<1/\rho.
\]
The positive root $1/\rho$ cannot lie on that circle, and there are no
other positive roots. Again Theorem \ref{thm:spectral} applies.

This argument permits multiple roots and roots on the unit circle.
No numerical information about the zeros of Eulerian polynomials is
required, and we have not used their stronger real-rootedness property.

\subsection{The two exceptional low exponents}

For $d=0$,
\[
A^\psi_{0,u}(z)=\frac{1-z}{1-(1+u)z},\qquad
a_0=1,\quad a_n=u(1+u)^{n-1}\quad(n\ge1).
\]
Thus $\Delt_1=-u$ and $\Delt_n=0$ for $n\ge2$.

For $d=1$,
\begin{equation}\label{eq:d1}
A^\psi_{1,u}(z)=\frac{(1-z)^2}{1-(u+2)z+z^2}
=1+\frac{uz}{1-(u+2)z+z^2}.
\end{equation}
The coefficients $a_n$ for $n\ge1$, extended by the auxiliary value
$\widetilde a_0=0$, satisfy
\[
\widetilde a_1=u,\qquad
\widetilde a_n=(u+2)\widetilde a_{n-1}-\widetilde a_{n-2}
\quad(n\ge2).
\]
For any sequence $x_{n+1}=h x_n-x_{n-1}$, the expression
$x_n^2-x_{n-1}x_{n+1}$ is invariant under shifting $n$: substitution
on both sides gives the same quadratic form.
Its value at the auxiliary index $1$ here is $u^2$. Hence the actual
$\Delt_n=u^2$ for $n\ge2$. At $n=1$, the actual $a_0=1$ instead gives
$\Delt_1=-2u$.
Together with the preceding two subsections, this proves
Theorem \ref{thm:power-main}.

\section{Square weights: exact recurrence, sharp gaps, and density}
\label{sec:squares}

We now fix $d=2$ and $u=1$. This is the original Challenge 2 sequence.
It is especially suitable for a detailed analysis because its generating
function has a cubic denominator.

\subsection{The known sequence and an indexing precaution}

From $E_2(z)=1+z$,
\begin{equation}\label{eq:square-gf}
A(z)=\frac{(1-z)^3}{1-4z+2z^2-z^3}
=1+\frac{z(1+z)}{1-4z+2z^2-z^3}.
\end{equation}
The coefficients begin
\[
(a_n)_{n\ge0}=1,1,5,18,63,221,776,2725,9569,33602,117995,\ldots.
\]
OEIS A033453 \cite{OEIS} starts at offset $0$ with
$1,5,18,63,221,\ldots$. Thus the precise relationship is
\begin{equation}\label{eq:OEIS-shift}
a_n=\mathrm{A033453}(n-1)\quad(n\ge1),\qquad a_0=1.
\end{equation}
The OEIS entry already gives the shifted generating function and the
recurrence with signature $(4,-2,1)$, as well as the interpretation by
compositions whose parts of size $k$ have $k^2$ colors.
These are established background facts, not claims of novelty here.

In our indexing,
\begin{equation}\label{eq:a-recurrence}
a_0=1,\ a_1=1,\ a_2=5,\ a_3=18,\qquad
a_n=4a_{n-1}-2a_{n-2}+a_{n-3}\quad(n\ge4).
\end{equation}
The starting index matters: applying the last recurrence at $n=3$ with
the actual $a_0=1$ would incorrectly give $19$.
The rational series $(z+z^2)/(1-4z+2z^2-z^3)$ has auxiliary coefficient
$0$ at $n=0$ and agrees with $a_n$ for $n\ge1$.
That auxiliary sequence satisfies the homogeneous recurrence from $n=3$.

\subsection{An exact determinant recurrence}

\begin{proposition}\label{prop:delta-recurrence}
The square-weighted determinant satisfies
\begin{align}
\Delt_1&=-4,&\Delt_2&=7,&\Delt_3&=9,&\Delt_4&=-9,\label{eq:delta-initial}\\
\Delt_n&=2\Delt_{n-1}-4\Delt_{n-2}+\Delt_{n-3}
&&&&(n\ge5).\label{eq:delta-recurrence}
\end{align}
Equivalently,
\begin{equation}\label{eq:delta-gf}
\sum_{n\ge2}\Delt_nz^{n-2}
=\frac{7-5z+z^2}{1-2z+4z^2-z^3}.
\end{equation}
Every $\Delt_n$ with $n\ge2$ is odd.
\end{proposition}

\begin{proof}
Let $\lambda,\beta,\overline\beta$ be the roots of
\begin{equation}\label{eq:f}
f(x)=x^3-4x^2+2x-1.
\end{equation}
The discriminant is $-107$, so the roots are distinct, with one real root
and one nonreal conjugate pair. Since $f(7/2)=-1/8$ and $f(4)=7$,
the real root satisfies $7/2<\lambda<4$.
Partial fractions in \eqref{eq:square-gf} give
\begin{equation}\label{eq:Binet}
a_n=\sum_{\xi\in\{\lambda,\beta,\overline\beta\}}c_\xi\xi^n,
\qquad c_\xi=\frac{\xi+1}{f'(\xi)},\qquad n\ge1.
\end{equation}
The same formula gives the auxiliary value $0$, not the actual $1$, at
$n=0$. None of its coefficients is zero.

For a sum $x_n=\sum_i c_i\xi_i^n$, pairwise expansion gives
\begin{equation}\label{eq:Cassini}
x_n^2-x_{n-1}x_{n+1}
=-\sum_{i<j}c_ic_j(\xi_i-\xi_j)^2(\xi_i\xi_j)^{n-1}.
\end{equation}
Indeed the diagonal terms cancel, while a pair contributes
$c_ic_j(\xi_i\xi_j)^{n-1}(2\xi_i\xi_j-\xi_i^2-\xi_j^2)$.
For the actual $a_n$, this formula is valid for $n\ge2$.

Vieta's formulas for \eqref{eq:f} give
\[
\lambda+\beta+\overline\beta=4,\quad
\lambda\beta+\lambda\overline\beta+\beta\overline\beta=2,\quad
\lambda\beta\overline\beta=1.
\]
The three pair-products in \eqref{eq:Cassini} therefore have elementary
symmetric functions $2,4,1$. Their polynomial is
$x^3-2x^2+4x-1$, proving \eqref{eq:delta-recurrence} for $n\ge5$.
Direct substitution of the initial $a_n$ gives \eqref{eq:delta-initial};
coefficient multiplication then gives \eqref{eq:delta-gf}.
Finally, reducing \eqref{eq:delta-recurrence} modulo $2$ gives
$\Delt_n\equiv\Delt_{n-3}\pmod2$. The three starting values
$\Delt_2,\Delt_3,\Delt_4$ are odd, proving the parity assertion.
\end{proof}

Some small exact values are useful both as a check and as a warning
against guessing eventual alternation.
\begin{center}
\begin{tabular}{rrrrrr}
\toprule
$n$ & $a_n$ & $\Delt_n$ & $n$ & $a_n$ & $\Delt_n$\\
\midrule
1&1&$-4$&7&2725&81\\
2&5&7&8&9569&311\\
3&18&9&9&33602&249\\
4&63&$-9$&10&117995&$-665$\\
5&221&$-47$&11&414345&$-2015$\\
6&776&$-49$&12&1454992&$-1121$\\
\bottomrule
\end{tabular}
\end{center}
In particular, equality is absent at every positive index, not merely
eventually.

\subsection{A single oscillatory term with a decaying positive correction}

Write $\beta=a+\ii b$ with $b>0$. The preceding identities yield
\begin{equation}\label{eq:beta-properties}
a=\frac{4-\lambda}{2}>0,\qquad
|\beta|^2=\beta\overline\beta=\lambda^{-1}.
\end{equation}
Put
\begin{align}
\rho&=\lambda^{-1},&\nu&=\lambda\beta,&|\nu|&=\sqrt\lambda,\label{eq:nu}\\
K&=-c_\lambda c_\beta(\lambda-\beta)^2,&
K_0&=4|c_\beta|^2b^2>0.\label{eq:Ks}
\end{align}
Here $K\ne0$ and $c_\lambda>0$.
Formula \eqref{eq:Cassini} becomes the exact expression
\begin{equation}\label{eq:delta-exact}
\Delt_n=2\Rea(K\nu^{n-1})+K_0\rho^{n-1}\qquad(n\ge2).
\end{equation}
Let $\theta=\arg\beta=\arg\nu$ and $\phi=\arg K$. Dividing by a
positive factor gives
\begin{equation}\label{eq:cosine}
\frac{\Delt_n}{2|K|\lambda^{(n-1)/2}}
=\cos\bigl((n-1)\theta+\phi\bigr)
 +\frac{K_0}{2|K|}\lambda^{-3(n-1)/2}.
\end{equation}
Thus the determinant is governed by a rotation angle, with an
exponentially decaying perturbation.

For orientation only, rounded numerical values are
\begin{align*}
\lambda&\approx3.511547141694532,&\rho&\approx0.284774761564910,\\
\beta&\approx0.244226429152734+0.474476778007327\ii,\\
\nu&\approx0.857612619217545+1.666147573612060\ii,\\
\theta&\approx1.095435949104736,&\theta/\pi&\approx0.348688092281161.
\end{align*}
The subsequent proofs use exact inequalities and algebraic identities,
not these decimal approximations.

\subsection{The sharp four-index block bound}

\begin{lemma}[Four equally spaced circle points]\label{lem:four-points}
Suppose $\pi/3<\theta<\pi/2$, and put
$\delta=-\cos(3\theta/2)>0$.
For every real $\alpha$, among
$\cos(\alpha+j\theta)$, $j=0,1,2,3$, there is a value at least $\delta$
and one at most $-\delta$.
\end{lemma}
\begin{proof}
The four circle points have circular gaps
$\theta,\theta,\theta,2\pi-3\theta$.
The largest gap is $2\pi-3\theta$ because $\theta<\pi/2$.
Every point of the circle is within angular distance at most
$\pi-3\theta/2$ of one of the four sample points.
Apply this to angles $0$ and $\pi$. Their nearest sample points have
cosines at least $\cos(\pi-3\theta/2)=\delta$ and at most $-\delta$,
respectively.
\end{proof}

\begin{proposition}\label{prop:four-block}
Every four consecutive indices $n\ge1$ contain both determinant signs.
The maximum possible length of a run of either sign is exactly three.
\end{proposition}
\begin{proof}
By \eqref{eq:beta-properties},
\begin{equation}\label{eq:cos-theta}
x:=\cos\theta=\frac{(4-\lambda)\sqrt\lambda}{2}.
\end{equation}
The function on the right is decreasing for $\lambda>4/3$.
Since $7/2<\lambda<4$, we obtain
$0<x<\sqrt{14}/8<1/2$, hence $\pi/3<\theta<\pi/2$.
For $0<x<1/2$ the function $4x^3-3x$ is decreasing. Therefore
\begin{equation}\label{eq:delta-margin}
\delta^2=\frac{1+\cos3\theta}{2}
=\frac{1+4x^3-3x}{2}
>\frac{64-17\sqrt{14}}{128}>\frac1{400}.
\end{equation}
The last inequality is exact: it is equivalent to
$17\sqrt{14}<1592/25$, which follows on squaring from
$4046<2534464/625$.
Thus $\delta>1/20$.

We also need a coarse bound on the correction in \eqref{eq:cosine}.
Because $f'(\lambda)=|\lambda-\beta|^2$ and
$|f'(\beta)|=2b|\lambda-\beta|$, the constants simplify to
\begin{equation}\label{eq:correction-ratio}
\frac{K_0}{2|K|}
=\frac{b|\beta+1|}{(\lambda+1)|\lambda-\beta|}
<\frac{|\beta|(1+|\beta|)}{(\lambda+1)(\lambda-|\beta|)}.
\end{equation}
Here $|\beta|=\lambda^{-1/2}<3/5$ and $\lambda>7/2$, so
\[
\frac{K_0}{2|K|}
<\frac{(3/5)(8/5)}{(9/2)(29/10)}
=\frac{32}{435}<\frac1{10}.
\]
For every $n\ge2$, the perturbation in \eqref{eq:cosine} is consequently
strictly less than $1/20$, since $\lambda^{-3/2}<1/2$.

On any block beginning at $n\ge2$, Lemma \ref{lem:four-points} gives
one cosine at least $\delta>1/20$ and one at most $-\delta<-1/20$.
The perturbations cannot destroy either corresponding sign.
The remaining block beginning at $n=1$ already contains
$\Delt_1=-4$ and $\Delt_2=7$.
Finally, $\Delt_4,\Delt_5,\Delt_6$ are negative, while
$\Delt_7,\Delt_8,\Delt_9$ are positive. Both possible run lengths of
three are attained, so the bound is sharp.
\end{proof}

This gives a particularly elementary proof of the infinitude assertion
in Challenge 2. The density conclusion requires one further fact:
the rotation angle is irrational relative to $\pi$.

\subsection{Certifying irrationality of the root angle}

\begin{lemma}\label{lem:irrational-angle}
For the roots of \eqref{eq:f}, $\theta/\pi$ is irrational.
\end{lemma}
\begin{proof}
The only possible rational roots of the monic cubic $f$ are $1$ and $-1$;
neither is a root. Hence $f$ is irreducible over $\Q$.
Its discriminant is
\[
(-4)^2 2^2-4\cdot2^3-4(-4)^3(-1)-27(-1)^2
+18(-4)\cdot2\cdot(-1)=-107.
\]
The splitting field $F$ therefore has Galois group $S_3$:
an irreducible cubic has a transitive subgroup of $S_3$, and a nonsquare
discriminant excludes its subgroup $A_3$.
The quadratic subfield fixed by $A_3$ is $\Q(\sqrt{-107})$.

Suppose $\theta/\pi$ were rational. Then
$\zeta=\beta/\overline\beta=e^{2\ii\theta}$ would be a root of unity
in $F$. The cyclotomic extension $\Q(\zeta)/\Q$ is Galois and abelian,
so the restriction homomorphism from $S_3$ to its Galois group kills
the commutator subgroup $A_3$. Thus
\[
\Q(\zeta)\subseteq\Q(\sqrt{-107}).
\]
A quadratic field can contain a root of unity of order $m$ only if
$\varphi(m)\le2$, giving $m=1,2,3,4,6$. The nonreal possibilities generate
$\Q(\sqrt{-3})$ or $\Q(\ii)$, neither of which equals
$\Q(\sqrt{-107})$. Hence $\zeta\in\{1,-1\}$.
But $\zeta=1$ would make $\beta$ real, while $\zeta=-1$ would make
$\beta+\overline\beta=0$, forcing $\lambda=4$ contrary to $f(4)=7$.
This contradiction proves the lemma.
\end{proof}

\subsection{Half-density in every arithmetic progression}

\begin{lemma}[Irrational rotation and a vanishing perturbation]
\label{lem:density}
Suppose $\theta/(2\pi)$ is irrational, $\phi\in\R$, and $\varepsilon_n\to0$.
Then the sets where
$\cos(n\theta+\phi)+\varepsilon_n$ is positive and negative each have
natural density $1/2$. The same relative-density conclusion holds after
restricting $n$ to any infinite arithmetic progression.
\end{lemma}
\begin{proof}
For every nonzero integer $k$,
\[
\frac1N\sum_{n=1}^N e^{\ii k(n\theta+\phi)}\longrightarrow0
\]
by the finite geometric-sum formula, since $e^{\ii k\theta}\ne1$.
Approximating continuous periodic functions by trigonometric polynomials,
and then approximating interval indicators from above and below, shows
that $n\theta+\phi$ is uniformly distributed modulo $2\pi$.

Fix $\eta>0$. Once $|\varepsilon_n|<\eta$, the perturbed sign agrees
with the cosine sign whenever $|\cos(n\theta+\phi)|>\eta$.
Uniform distribution gives the density of this exceptional strip as its
normalized angular length, which tends to zero as $\eta\downarrow0$.
The positive and negative cosine regions each have length $\pi$.
Squeezing between the regions with thresholds $\eta$ and $-\eta$ proves
the two densities.

On the progression $n=r+jm$, replace $\theta$ by $m\theta$ and the phase
by $r\theta+\phi$. Irrationality persists, and the restricted perturbation
still tends to zero. The same proof applies.
\end{proof}

Apply Lemma \ref{lem:density} to \eqref{eq:cosine}, using
Lemma \ref{lem:irrational-angle}. This proves
\eqref{eq:AP-density} and the corresponding negative-density formula.
Together with Propositions \ref{prop:delta-recurrence} and
\ref{prop:four-block}, it proves all three parts of
Theorem \ref{thm:square-main}.
If the sign sequence had an eventual period $m$, it would be eventually
constant on each residue class modulo $m$, contradicting density $1/2$
for both signs within that class.

Notice what has not been needed: no Diophantine lower bound on how closely
the angle can approach a zero of the cosine is required. The density
argument discards a shrinking-boundary issue using fixed strips and then
lets the strip width tend to zero. It does not assert that the cosine
and determinant signs agree at every sufficiently large individual index.

\section{Exact computation and reproducibility}
\label{sec:computation}

The program \texttt{verify.py} uses only the Python standard library
and runs with Python 3.9 or later. It makes no network calls and uses
integers or rational numbers for every asserted coefficient, determinant,
or certified sign. Its default run was executed with Python 3.13.5.
The output is included in \texttt{verification-output.txt} and in the
\texttt{data/} directory.

\subsection{Independent checks rather than one implementation twice}

The square coefficients are generated in two ways: by the convolution
\eqref{eq:convolution} with weights $k^2$, and by coefficient extraction
from the rational function \eqref{eq:square-gf}. They agree through
$a_{601}$. Direct adjacent determinants from those coefficients agree
with the independent determinant recurrence \eqref{eq:delta-recurrence}
through $\Delt_{600}$.

The same convolution-versus-rational comparison is performed for
$d=0,\ldots,12$ and $u=1,2,3$, through coefficient $a_{121}$.
Eulerian symmetry, positivity, $E_d(1)=d!$, and the relevant denominator
identities are checked through $d=30$.
For divisor weights, a sieve is compared with direct divisor enumeration,
and \eqref{eq:sigma-even} is checked for $d=0,\ldots,10$ and
$n=1,\ldots,300$.
These finite checks are useful for catching normalization and indexing
errors, especially the exceptional actual value $a_0=1$ in the square
and linear-power formulas.

The long square computation uses only three live determinant integers,
updated by \eqref{eq:delta-recurrence}. This avoids computing a small
relative difference of two enormous, nearly equal floating-point products.
The integers grow exponentially in magnitude, so constant \emph{number}
of live integers does not mean constant bit-space. The algorithm uses
$O(N)$ bits of working integer storage up to index $N$, excluding saved
output, and $O(N)$ recurrence steps with increasing operand lengths.

\begin{center}
\begin{tabular}{rrrrr}
\toprule
Range $1\le n\le N$ & $\Delt_n>0$ & $\Delt_n<0$ & $\Delt_n=0$ & Positive fraction\\
\midrule
100&50&50&0&0.500000\\
1,000&501&499&0&0.501000\\
10,000&5,003&4,997&0&0.500300\\
100,000&50,000&50,000&0&0.500000\\
\bottomrule
\end{tabular}
\end{center}
The exact equality of the last two counts is a property of this finite
cutoff, not a claimed formula at every cutoff. The universal density
statement is established by the irrational-rotation proof.
The longest observed runs have length three, agreeing with the proven
sharp four-index block bound. The archive also contains counts in every
residue class for moduli $1$ through $12$.

\subsection{Rational certificates for negative Lambert values}

For sample nonnegative integer exponents, the program certifies
$L_d(-9/10)>1$ without trusting floating-point evaluation of an alternating
series. Let $0<r<1$, let
\[
S_N=\sum_{n=1}^N\sig_d(n)(-r)^n,
\qquad q=r\left(\frac{N+2}{N+1}\right)^{d+1}<1.
\]
Since $\sig_d(n)\le n^{d+1}$ and the ratios of successive majorant terms
$n^{d+1}r^n$ decrease with $n$, the remainder obeys
\begin{equation}\label{eq:tail-certificate}
|L_d(-r)-S_N|
\le\frac{(N+1)^{d+1}r^{N+1}}{1-q}=:T_N.
\end{equation}
For rational $r$, both $S_N$ and $T_N$ are exact rational numbers.
The assertion $S_N-T_N>1$ is then a finite, exact certificate of a
negative-side crossing of $1$.

With $r=9/10$ and $N=600$, this succeeds for all $d=0,\ldots,10$.
For example, the following numbers are rounded displays of the exact
certificates saved in the JSON file:
\begin{center}
\begin{tabular}{rrr}
\toprule
$d$ & $S_{600}$, approximately & $T_{600}$, approximately\\
\midrule
0 & $7.09052786658$ & $1.93\times10^{-24}$\\
1 & $32.3413228274$ & $1.18\times10^{-21}$\\
2 & $255.356970620$ & $7.18\times10^{-19}$\\
5 & $1.39444495497\times10^6$ & $1.64\times10^{-10}$\\
10 & $9.98138686366\times10^{13}$ & $1.40\times10^4$\\
\bottomrule
\end{tabular}
\end{center}
Even the coarse final absolute error is negligible relative to that
partial sum and still gives a rigorous lower bound greater than $1$.
These certificates are independent numerical corroboration of the
analytic pole mechanism. The proof for arbitrary real $s$ and positive
$u$ uses \eqref{eq:negative-radial}, not a finite grid of examples.

\subsection{Rebuilding the archive contents}

From the unpacked directory, run
\begin{verbatim}
python verify.py
pdflatex -interaction=nonstopmode -halt-on-error article.tex
pdflatex -interaction=nonstopmode -halt-on-error article.tex
\end{verbatim}
The first command regenerates the CSV, JSON, and sign data. The two LaTeX
runs resolve the table of contents and cross-references. A standard TeX
Live installation with the packages listed in the preamble suffices;
no external fonts are bundled. A shorter exploratory run is available
through \texttt{python verify.py --limit 10000}.
An optional \texttt{numeric\_square.py} script uses \texttt{mpmath} to
reproduce the rounded root and amplitude values; it is not needed for
any exact check or proof.

\section{Scope, priority, and remaining questions}
\label{sec:scope}

\subsection{What has been established}

The original two targets have affirmative proofs. In the divisor family,
the argument establishes both strict signs, with bounded gaps, for
all real exponents and all positive composition weights. In the power
family it classifies every nonnegative integer exponent, again for every
positive weight. In the unweighted square case, the exact recurrence,
nonvanishing, optimal four-index block bound, and half-density on every
arithmetic progression give a substantially more explicit description
than infinitude alone.

The general pole theorem is deliberately formulated to survive repeated
secondary poles, ties in modulus, and rational dependencies between
angles. The square density argument is more specialized: it uses the
single conjugate pair and a verified irrational angle. One should not
transfer its precise density $1/2$ to every parameter in the larger
families merely from the bounded-gap theorem.

\subsection{What is background, and what priority is being claimed}

The geometric-transform definition, colored-composition interpretation,
Eulerian generating functions, partial fractions, Pringsheim's principle,
and irrational-rotation equidistribution are standard tools or are
explicitly present in the cited sources. They are not presented as new
mathematical concepts. The A033453 generating function and recurrence are
already recorded in the OEIS. The contribution developed here is the
proof of the specified challenges and the accompanying stronger sign
statements and classification.

The bibliographical audit was focused rather than exhaustive. The arXiv
record inspected lists the 2023 first version; its Section 3 expressly
poses the two questions. The publisher metadata identifies a journal
publication with DOI \texttt{10.1142/S1793042124500192}; the revised
publisher full text was not available for direct comparison in this
session. The challenge numbering in this manuscript therefore refers
unambiguously to the inspected arXiv version, not to an uninspected
revision. Searches for the paper, its identifier, its authors with
geometric log-concavity, and A033453 did not surface a prior proof of
these specific assertions. Search coverage was uneven, so this absence
is evidence for choosing the problem, not a proof of originality.

A 2026 paper by Alzer and Volkmer \cite{AV26} was also inspected because
its title combines Lambert series and inequalities. Its main
log-concavity statement concerns the real-variable function
$t\mapsto L_0(1-e^{-t})$, not the coefficients of $1/(1-uL_s(z))$.
That functional inequality is a different question and does not supply
the coefficient-sign conclusions proved here.

\subsection{Questions naturally left by the proof}

For small nonnegative integer divisor exponents, the computational
odd/even pattern is much stronger than mere infinitude. Establishing
such an exact parity pattern, where it is valid, would require more
specific information about the competing poles or a different direct
inequality. The present proof deliberately avoids assuming it.

For the larger parameter families, explicit optimal gap bounds and exact
sign densities would require understanding the first subdominant spectrum.
Theorem \ref{thm:spectral} provides a route to a computable bound from
certified pole locations and residues, but does not compute such constants
uniformly in the parameters. Changes in which poles have smallest
nondominant modulus can create transitions in the oscillatory pattern.

Finally, higher-order adjacent Hankel or Tur\'an determinants provide a
natural extension. Their leading geometric contributions cancel even
more extensively. A corresponding analysis would have to identify the
first nonvanishing combination of several pole contributions, rather
than simply reuse the two-term cross expression \eqref{eq:delta-cross}.
These are further problems, not premises needed by the results above.

\appendix
\section{Proof dependencies and verification boundaries}

The main analytic route is
\[
\text{positive weights}
\ \Longrightarrow\ \text{unique simple positive dominant pole}
\ \Longrightarrow\ \text{subdominant analysis}.
\]
For divisor weights, the missing nondominant pole is supplied by
\[
\text{divisor parity identity}
\ \Longrightarrow\ L_s(-r)\to+\infty
\ \Longrightarrow\ \text{a negative pole inside }|z|<1.
\]
For power weights, it is supplied by
\[
\text{Eulerian rational function}
\ \Longrightarrow\ \text{sign or reciprocal-root analysis}
\ \Longrightarrow\ \text{a nonpositive first subdominant circle}.
\]
The last phrase means that the circle contains \emph{no pole on the
positive real axis}; it does not mean every pole is a nonpositive real number.
The common spectral theorem then proves bounded gaps for both signs.

The square-specific route is algebraic:
\[
\text{cubic denominator}
\ \Longrightarrow\ \text{pair-product determinant recurrence}
\ \Longrightarrow\ \text{parity and exact cosine formula}.
\]
Root bounds and four-point circle geometry prove the sharp gap bound.
The Galois group $S_3$ proves irrationality of the angle, and irrational
rotation proves the density theorem.

All numerical roots in the text are illustrative approximations.
No root isolation obtained numerically is substituted for the exact
inequalities $7/2<\lambda<4$, the discriminant calculation, or the
Galois argument. The Python tests do not constitute a formalization of
complex analysis or Galois theory; they are independent exact checks of
the algebraic identities and finite data used to audit the manuscript.

\begin{thebibliography}{9}

\bibitem{HN23}
B.~Heim and M.~Neuhauser,
\emph{Log-Concavity of Infinite Product and Infinite Sum Generating Functions},
arXiv:2302.13327v1, 26 February 2023.
Challenges 2 and 3 occur in Section 3, printed page 10.
\url{https://arxiv.org/abs/2302.13327}.
Journal publication DOI: \url{https://doi.org/10.1142/S1793042124500192}.
The proof targets are cited to the inspected arXiv version.

\bibitem{HN21}
B.~Heim and M.~Neuhauser,
\emph{Horizontal and vertical log-concavity},
Research in Number Theory \textbf{7}, article 18 (2021).
\url{https://doi.org/10.1007/s40993-021-00245-1}.

\bibitem{OEIS}
The OEIS Foundation Inc.,
\emph{The On-Line Encyclopedia of Integer Sequences}, entry A033453,
``INVERT transform of squares.''
\url{https://oeis.org/A033453}.
Accessed 19 September 2026. The entry's offset is $0$ and its first
term corresponds to $a_1$ in this manuscript.

\bibitem{AV26}
H.~Alzer and H.~W.~Volkmer,
\emph{Inequalities for the Lambert series},
The Ramanujan Journal \textbf{70}, article 1 (2026), published 10 April 2026.
\url{https://doi.org/10.1007/s11139-026-01365-x}.
Cited to distinguish functional log-concavity from the coefficient
problem studied here.

\end{thebibliography}
\end{document}
